%% LyX 2.1.4 created this file.  For more info, see http://www.lyx.org/.
%% Do not edit unless you really know what you are doing.
\documentclass[unicode,master,pdfcover]{main}
\usepackage{fontspec}
\usepackage{color}
\usepackage[unicode=true,
 bookmarks=true,bookmarksnumbered=true,bookmarksopen=false,
 breaklinks=false,pdfborder={0 0 1},backref=false,colorlinks=false]
 {hyperref}
% table
\usepackage{xspace}
% 告诉 xspace 遇到这些中文标点时不要加空格
\xspaceaddexceptions{（ ） ， 。 ” ’ ！ ？ ： ；}
\usepackage{amsfonts}
\usepackage{algorithmic}
\usepackage{placeins}
% \usepackage{stfloats}
% % \usepackage{url}
% \usepackage{verbatim}
% \usepackage[table]{xcolor}

% \usepackage{paralist}
\usepackage{multicol}
\usepackage{multirow}
% \usepackage{subcaption}
\usepackage{makecell}
% \usepackage{rotating}
% \usepackage{caption}
\usepackage{subcaption}
\usepackage{minted}
\usepackage{pdfpages}

% \usepackage{etoolbox}
% \usepackage{stfloats}
% \usepackage{mathtools}
\usepackage{bbding}

\addbibresource{main.bib}

% \setlength{\bibsep}{0pt plus 2pt minus 1pt}%参考文献间距
% 定义一个新的正文引用命令
% \newcommand{\mycite}[1]{{\setcitestyle{numbers,square,sort&compress}\cite{#1}}}

% 全局将名字改为中文“算法”
\floatname{algorithm}{算法}
% \renewcommand{\thealgorithm}{\thechapter-\arabic{algorithm}}

\newcommand{\theHalgorithm}{\arabic{algorithm}}
\newcommand{\vect}[1]{\mathbf{#1}} % 向量/矩阵：加粗正体
\newcommand{\matr}[1]{\mathbf{#1}} % 备用名：矩阵
\newcommand{\termfix}{%
  \ifhmode
    \ifdim\lastskip>0pt
      \unskip
      % Keep one interword space after punctuation (e.g., comma) in English text.
      \ifnum\spacefactor>1000\space\fi
    \fi
  \fi
}
% \newcommand{\sexyname}{SG-Mamba\xspace}
\newcommand{\sexyname}{\termfix AC-Mamba\xspace}
\newcommand{\fullnameen}{\termfix Anatomy-Calibrated Mamba for Temporal Micro-Action Localization\xspace}
% \newcommand{\fullname}{Skeleton-Guided Mamba for Temporal Micro-Action Localization\xspace}
\newcommand{\fullname}{\termfix 基于解剖结构校准 Mamba\xspace}
\newcommand{\crossblockfull}{\termfix 解剖结构校准Mamba\xspace}
\newcommand{\channelblockfull}{\termfix 冗余剔除层\xspace}
\newcommand{\crossblockfullen}{\termfix Anatomy-Calibrated Mamba\xspace}
\newcommand{\channelblockfullen}{\termfix Redundancy Removal Layer\xspace}
\newcommand{\crossblockshort}{\termfix ACM\xspace}
\newcommand{\channelblockshort}{\termfix RRL\xspace}
\newcommand{\featuremodulefull}{\termfix 多尺度动态特征模块\xspace}
\newcommand{\featuremodulefullen}{\termfix Multi-Scale Dynamic Feature Module\xspace}
\newcommand{\featuremoduleshort}{\termfix MSDF\xspace}
\newcommand{\fulltwo}{\termfix 时序均衡与逻辑演化的青光眼手术动作检测框架\xspace}
\newcommand{\fulltwoen}{\termfix Temporal-Balance and Logic-Evolution Network for Glaucoma Surgical Online Action Detection\xspace}
\newcommand{\sexytwo}{\termfix TBL Net\xspace}
\newcommand{\graphmodulefull}{\termfix 规程引导模块\xspace}
\newcommand{\graphmodulefullen}{\termfix Procedural Guidance Constrainer\xspace} 
% Adaptive Protocol Evolution Module / Logic Evolution Augmentor
\newcommand{\graphmoduleshort}{\termfix PGC\xspace}
\newcommand{\lossfull}{\termfix 类别边界加权损失函数\xspace}
\newcommand{\lossshort}{\termfix CBD Loss\xspace}
\newcommand{\lshort}{\termfix CBD\xspace}
\newcommand{\lossfullen}{\termfix Category-Boundary Dual-Weighted Loss\xspace} 

\newcommand{\memoryunit}{抗漂移记忆单元}
\newcommand{\evolution}{时长感知演化单元}
\newcommand{\fusion}{动态门控融合单元}
% Duration-Aware Boundary Enhancement Loss
\definecolor{mygray}{rgb}{0.9,0.9,0.9}
\def\etal{\emph{et~al}.}
% \newcommand{\cell}{\cellcolor[rgb]{0.92,0.92,0.92}}
\hypersetup{
  pdftitle={医疗场景下的视频动作定位与识别研究},
  pdfauthor={李奕锐},
  pdfsubject={华南理工大学专业硕士学位论文},
  pdfkeywords={关键字1, 关键字2},
  unicode=true,linkcolor=black, anchorcolor=black, citecolor=black, filecolor=magenta, menucolor=red, urlcolor=black, pdfstartview=FitH}
\makeatletter

%%%%%%%%%%%%%%%%%%%%%%%%%%%%%% LyX specific LaTeX commands.
\providecommand{\LyX}{\texorpdfstring%
  {L\kern-.1667em\lower.25em\hbox{Y}\kern-.125emX\@}
  {LyX}}
%% Because html converters don't know tabularnewline
\providecommand{\tabularnewline}{\\}

\makeatother

\usepackage{xunicode}
\renewcommand{\lstlistingname}{列表}

\begin{document}

\title{医疗场景下的视频\\动作定位与识别研究}
\author{李奕锐}
\supervisor{指导教师：陈健\ 教授}
\institute{华南理工大学}


\date{2026年3月8日}

\maketitle
\frontmatter

\begin{abstractCN}
% 硕士学位论文的中文摘要一般约500～800字。摘要内容应涉及本项科研工作的目的和意义、研究思想和方法、研究成果和结论
近年来，深度学习在诸多医疗影像分析任务中取得了长足进步，
并在手术阶段识别、解剖结构追踪等任务上取得了突破性成果。
然而，深度学习在医疗视觉任务上的成功往往建立在大规模、高质量的专家标注数据集基础之上。
在临床实际应用中，医疗视频数据的构建常面临影像对比度低、动作幅度微小、数据分布极度失衡等困难，
使得深度模型难以学习到有效的视觉表征以完成精准的动作理解任务。
因此，如何深度融合医学领域先验，实现更高效的医疗视频动作表征学习，对推动智慧医疗系统的实际落地具有重要意义。
然而，针对辅助诊断的吞咽造影视频和面向术中导航的青光眼手术视频两种典型视觉数据，高效的表征学习仍存在以下挑战：
（1）在吞咽造影视频微动作定位中，X 射线影像的低对比度及关键结构的模糊遮挡，导致微弱动作信号极易被静态背景噪声淹没，使得模型难以精准捕捉时序边界；
（2）在眼科手术视频在线识别中，手术阶段时长的显著失衡以及不同操作间极高的外观相似性，导致模型在实时推理时极易产生漏检与误判。
针对上述挑战，本文研究内容如下：

对于吞咽造影视频分析，本文针对时序微动作定位任务中解剖结构模糊及弱目标特征淹没的问题，对解剖结构校准机制进行探索。
提出了一种名为 \sexyname 的全新框架，通过引入基于临床先验的骨架热图模态，为模型提供鲁棒的几何结构参照，
从而引导模型学习抗干扰能力更强的视觉表征。
此外，为保证多模态特征的高效校准与增强，该方法利用跨模态特征校准模块 \crossblockshort 和通道级增强机制 \channelblockshort，
通过骨架特征引导空间聚焦并显式剥离冗余的静态背景信号，实现对细微动作信号的有效放大。
在 VFSS 基准数据集上的实验表明，相较于现有方法，所提出的方法在平均 mAP 指标上取得了 5.3\% 的性能提升，并在外部验证集中证实了其卓越的泛化能力。

% 这里只谈了显微手术，但是没有聚焦青光眼
对于青光眼手术导航，本文针对在线手术动作识别任务中时序逻辑混淆及数据分布不均的问题，
对手术规程逻辑协同机制进行探索，提出了一种名为 \sexytwo 的动作检测框架。该方法旨在利用手术流程的结构化先验，
通过动态逻辑演化模块 \graphmoduleshort，提升模型对手术状态跳转及动作持续时长的感知能力，
从而引导模型学习逻辑连贯性更强的场景表征。此外，还提出了类别边界加权损失 \lossshort，以促进长短时动作特征的均衡学习。
在构建的 MIGS 数据集及公开数据集 OphNet 上的实验表明，相较于现有方法，本文所提出的方法分别取得了 0.9\% 和 1.9\% 的在线识别准确率提升。
\end{abstractCN}
\keywordsCN{医疗视频理解；吞咽造影视频；时序微动作定位；手术视频；在线动作识别}
\begin{abstractEN}
In recent years, deep learning has achieved remarkable progress in video action understanding and analysis tasks within the general domain.
However, a pronounced domain gap persists when migrating such technologies to the healthcare sector.
In clinical practice, compared with natural video analysis, medical video analysis is frequently plagued by low image contrast, subtle motion amplitudes, and highly imbalanced data distributions,
which make it difficult for deep models to learn effective visual representations for accurate action analysis.
Therefore, exploring how to deeply integrate medical domain priors to realize more efficient medical video action understanding is of great significance for the practical deployment of smart healthcare systems.
Focusing on two typical types of medical video data,
namely videofluoroscopic swallowing study (VFSS) videos for auxiliary diagnosis and glaucoma surgery videos for intraoperative navigation,
this paper thoroughly investigates the core challenges of medical video action understanding in complex clinical scenarios as follows:
(1) For temporal micro-action localization in VFSS videos,
the low contrast of X-ray images and the blurring and occlusion of key anatomical structures
easily cause weak motion signals to be overwhelmed by static background noise,
making it hard for models to capture precise temporal boundaries;
(2) For online action recognition in glaucoma surgery videos,
the significant imbalance in the duration of surgical phases and the high visual similarity among different surgical operations
easily lead to missed detections and false judgments during real-time inference.

To address the above challenges, the research contents of this dissertation are presented as follows:

For VFSS video analysis,
this dissertation proposes a novel framework named \sexyname
to tackle the problems of blurred anatomical structures and obscured weak target features in temporal micro-action localization tasks.
This framework introduces a skeleton heatmap modality based on clinical priors to provide the model with robust geometric structural references,
thus guiding it to learn more interference-resistant visual representations.
Furthermore, to ensure efficient calibration and enhancement of multimodal features,
the proposed method employs a Cross-modal Feature Calibration module (\crossblockshort)
and a Channel-level Enhancement mechanism (\channelblockshort).
These modules guide spatial focus through skeleton features and explicitly strip redundant static background signals,
thereby achieving effective amplification of subtle motion signals.
Experimental results on the VFSS benchmark dataset demonstrate that the proposed method
improves the mean Average Precision (mAP) by 5.3\% compared with state-of-the-art methods,
and it exhibits excellent generalization performance on an external validation set.

For glaucoma intraoperative navigation,
this dissertation proposes an online action detection framework named \sexytwo
to address the issues of temporal logic confusion and imbalanced data distribution in online surgical action recognition tasks.
By leveraging the structural priors of surgical procedures through a Dynamic Logic Evolution module (\graphmoduleshort),
this method enhances the model's ability to perceive surgical state transitions and action durations,
thus guiding it to learn more logically coherent scene representations.
Additionally, a Class Boundary Weighted Loss (\lossshort) is proposed to realize balanced learning of features for long and short-duration actions.
Extensive experiments on the self-constructed MIGS dataset and the public OphNet dataset
show that the proposed method achieves an improvement of 0.9\% and 1.9\% in online recognition accuracy,
compared with state-of-the-art methods on the two datasets respectively.

\end{abstractEN}

\keywordsEN{Medical Video Understanding; Videofluoroscopic Swallowing Study; Temporal Micro-action Localization; Glaucoma Surgery Video; Online Surgical Action Detection}

\tableofcontents{}

\listoftables
\listoffigures

% \include{sec/symbol}
\chapterx{英文缩略词}

\begin{table}
\centering{}%
\begin{tabular}{ccc}
VFSS & Video Fluoroscopic Swallowing Study & 视频透视吞咽检查 \tabularnewline
TMAL & Temporal Micro-Action Localization & 时序微动作定位 \tabularnewline
TAL & Temporal Action Localization & 时序动作定位 \tabularnewline
SSM & State Space Model & 状态空间模型 \tabularnewline
\crossblockshort & \crossblockfullen & \crossblockfull \tabularnewline
\channelblockshort & \channelblockfullen & \channelblockfull \tabularnewline
mAP & mean Average Precision & 平均精度均值 \tabularnewline
tIoU & temporal Intersection over Union & 时序交并比 \tabularnewline
PCK & Probability of Correct Keypoint & 关键点正确概率 \tabularnewline
FLOPs & Floating Point Operations & 浮点运算次数 \tabularnewline
CNN & Convolutional Neural Network & 卷积神经网络 \tabularnewline
RNN & Recurrent Neural Network & 循环神经网络 \tabularnewline
RGB & Red, Green, Blue & 红绿蓝（外观模态） \tabularnewline
EMA & Exponential Moving Average & 指数移动平均 \tabularnewline
OAD & Online Action Detection & 在线动作检测 \tabularnewline
SPR & Surgical Phase Recognition & 手术阶段识别 \tabularnewline
MIGS & Minimally Invasive Glaucoma Surgery & 微创青光眼手术 \tabularnewline
GCN & Graph Convolutional Network & 图卷积网络 \tabularnewline
GRU & Gated Recurrent Unit & 门控循环单元 \tabularnewline
MLP & Multi-Layer Perceptron & 多层感知机 \tabularnewline
ZOH & Zero-Order Hold & 零阶保持 \tabularnewline
\graphmoduleshort & \graphmodulefullen & \graphmodulefull \tabularnewline
% \featuremoduleshort & \featuremodulefullen & \featuremodulefull \tabularnewline
\lossshort & \lossfullen & \lossfull \tabularnewline
\end{tabular}
\end{table}
\mainmatter

\include{sec/1-intro}
\chapter{相关研究综述}

\section{时序（微）动作定位分析}
\subsection{通用时序动作定位方法}
时序动作定位（Temporal Action Localization, TAL）是视频理解领域的核心任务之一，
旨在从未经剪辑的长视频中识别出所有动作实例的类别，并精确回归其起始与终止时间。
本章从架构设计、特征提取范式以及序列建模机制三个维度对现有的TAL方法进行系统性的综述与分析。

\textbf{架构设计演进}\quad
根据网络架构设计的不同，现有的主流方法通常被划分为三大类~\cite{liu2024end}：
两阶段\cite{chao2018rethinking,lin2019bmn,zeng2019graph,zhao2021video,zhu2021enriching}、
一阶段\cite{zhang2022actionformer, shi2023tridet, yang2024adapting}
以及基于查询\cite{liu2022end}的方法。
早期的两阶段方法遵循“先生成后分类”的范式。
这类方法首先通过滑动窗口或边界采样生成一系列动作候选片段，
随后利用特定的采样模块提取特征进行分类与精修~\cite{chao2018rethinking, lin2019bmn,zhu2021enriching}。
典型的代表如BMN~\cite{lin2019bmn}通过设计双分支网络分别处理起始与终止边界的概率估计，在边界质量上取得了显著提升。
尽管两阶段方法在边界质量上具有优势，但其串行的处理流程限制了推理效率。
相比之下，一阶段方法直接基于多尺度特征金字塔进行端到端的动作定位，极大地简化了检测流程。
以ActionFormer\cite{zhang2022actionformer}为代表的先进模型，
利用Transformer\cite{vaswani2017attention}强大的上下文建模能力，
在单一网络中并行完成动作分类与边界回归，成为了当前性能与效率平衡的标杆。
此外，受目标检测DETR\cite{zhu2020deformable}架构启发，
基于查询的方法（如TadTR~\cite{liu2022end}）通过学习一组稀疏的查询向量与视频特征交互，
实现了无需后处理的直接预测。

\textbf{特征提取范式：从特征基到端到端}\quad
除了架构差异，TAD方法还可根据输入模态的处理方式分为“基于特征”和“端到端”两类。
目前的主流方法（如 ActionFormer）大多属于前者，
即依赖于预训练的视频编码器（如 I3D\cite{carreira2017quo}, VideoMAE\cite{tong2022videomae,wang2023videomae}）离线提取 RGB 或光流特征。
这种解耦策略虽然降低了显存开销，但忽略了特征提取与定位任务的协同优化。
近年来，端到端方法逐渐受到关注。
虽然早期工作如AFSD~\cite{lin2021learning}受限于计算资源需对输入进行大幅下采样，
但后续工作如的DaoTAD~\cite{wang2021rgb}和E2E-TAD~\cite{liu2022empirical}证明了仅凭RGB模态结合强数据增强即可实现优异性能。
为进一步解决显存瓶颈，SGP~\cite{cheng2022stochastic} 和 ETAD~\cite{liu2023etad}提出了部分反向传播策略，
而Re2TAL~\cite{zhao2023re2tal} 则设计了可逆网络架构以释放中间激活占用的显存。
AdaTAD~\cite{liu2024end}进一步将模型的尺寸扩大到10亿参数级别，展示了端到端方法在大规模模型上的潜力。
尽管端到端方法发展迅速，但在考虑模型显存和计算量的情况下，基于特征的方法依旧占据着重要地位。

\textbf{序列建模机制的推演}\quad
时序建模机制经历了从局部感知到全局关联，再到动态选择的选择性演进过程。
早期方法主要依赖一维卷积神经网络（1D-CNN）提取时序特征~\cite{shou2016temporal,lin2017single, xu2020g}，
虽具备局部归纳偏置，但为了扩大感受野而进行的下采样操作往往导致微小动作的细微特征丢失；
随后引入的图卷积网络（GCN）\cite{zeng2019graph,zhao2021video} 尝试将视频片段建模为图节点以聚合非局部语义关系。
为了获取全局上下文，基于Transformer的方法利用自注意力机制~\cite{zhang2022actionformer, yang2024adapting}实现了对长距离依赖的建模，
确立了当前的主流范式，但其二次的计算复杂度限制了高分辨率序列的解析，且静态注意力权重容易被动态的背景“劫持”，难以聚焦微弱的瞬态运动信号。
以 Mamba 为代表的状态空间模型（SSMs）\cite{gu2023mamba, li2024videomamba}
凭借线性复杂度与输入依赖的选择性扫描机制（Selective Scan）建模得到更精细的特征。
不同于传统的静态加权，这种机制赋予了模型内容感知的动态滤波能力，使其能够自适应地过滤冗余的静态背景并敏锐响应微弱的动态变化，
为在强背景干扰下实现精准动作定位提供了全新的技术路径。

\subsection{微动作时序定位研究}



\section{在线动作检测}


\section{医疗场景下的视频动作理解}
\subsection{吞咽造影视频分析}
吞咽造影（Videofluoroscopic Swallowing Study, VFSS）是临床评估吞咽功能的“金标准”，
其核心在于通过分析吞咽过程中的时序逻辑与运动特征来辅助诊断。
早期的研究主要聚焦于半自动分析技术，通过人工干预与计算机算法的结合来提取空间运动学参数。
例如，Kellen等\cite{kellen2010computer}与Hossain等\cite{hossain2014semi}先后探索了基于模板匹配与边缘特征追踪的方法，
利用颈椎等解剖学标志构建相对坐标系，
从而实现舌骨位移、轨迹及活动范围（ROM）的量化。
尽管此类方法在一定程度上提升了客观性并降低了分析耗时，但其高度依赖人工初始化，
且对低对比度或存在解剖重叠的造影图像表现出较强的敏感性，难以实现大规模临床数据的自动化处理。

随着深度学习技术的兴起，自动化时间学参数分析成为了该领域的研究热点方向之一，
其目标是自动识别吞咽各阶段（如口腔期、咽期、食管期）并提取关键时间参数（如咽传输时间、上食管括约肌开合时间等）。
早期的自动化尝试多基于卷积神经网络（CNN）开展。
研究者利用2D或3D CNN对视频帧进行特征提取，通过二分类任务区分咽期与非咽期\cite{lee2019automatic,bandini2021effect,jeong2023application}。
例如，Lee等\cite{lee2019automatic}利用预训练的3D CNN对咽期进行检测，验证了深度学习在VFSS视频分类中的可行性；
Jeong等\cite{jeong2023application}则进一步提出了“从分类到定位”的框架，实现了对剪辑视频中七个不同吞咽阶段的同时识别。
然而，临床实际获取的VFSS数据通常为未剪辑的长视频，且包含大量非典型或微小的动作变化，
这使得简单的分类模型难以满足高精度的临床需求。
针对这一挑战，Ruan等\cite{ruan2023temporal}针对未剪辑视频，
采用多检测器协同的“由粗到细”定位策略，实现了对微吞咽动作的识别与分类。
尽管目前的基于视觉流（RGB）或光流的方法取得了一定进展，但在VFSS场景下仍面临组织对比度低、背景噪声干扰严重等问题。
\subsection{手术工作流识别}

\section{本章小结}
\include{sec/3-alg1}
\chapter{基于的青光眼手术视频导航}

\section{引言}
青光眼（Glaucoma）作为全球范围内导致不可逆性失明的首要病因，
已成为严重的公共卫生负担\cite{flaxman2017global,sabanayagam2017global}。
微创青光眼手术凭借其创伤小、术后恢复快等优势，
已成为当前临床治疗的重要手段\cite{song2022minimally,zhang2023influence}。
然而，在眼科手术领域，手术空间极度狭小且解剖结构复杂。
前房角区域包含多层半透明且细微的生理结构，结构错综复杂，不同个体形态差异大，
即视为非青光眼的眼科专家，也同样难以快速识别\cite{dietlein2000morphological}。
统计数据显示，培养一名熟练的眼科医生需要漫长的学徒制教学周期，
而我国每 10 万人中眼科医生仅有不到4名，比例远低于发达国家\cite{nhc2022eyehealth, resnikoff2020estimated}，
优质医疗资源的紧缺与陡峭的手术学习曲线使得医疗事故风险依然存在。

在手术操作过程中，外科医生需要在复杂的动态场景下做出瞬时决策。
任何细微的识别延迟或步骤误判，都可能引发前房积血或眼压波动等严重并发症。
因此，实现在线手术阶段识别具有至关重要的临床意义。
在线手术阶段识别是构建“术中主动安全预警系统”的核心，
它不仅能为初级医生提供实时的规程指引，帮助其建立正确的认知框架，
还能作为下游任务（如风险预测、自动器械追踪）的基础，
从而真正将人工智能技术整合进手术流程，提升操作的标准化程度。

% 尽管针对普通外科等大尺度动作手术（如胸腹腔手术）的识别研究已取得一定进展
% \cite{endonet, czempiel2020tecno, wang2022autolaparo, lavanchy2024challenges, liu2025lovit}，
% 但
在微创青光眼手术这一高度精细且实时性要求极高的场景下，实现鲁棒的在线识别仍面临以下核心挑战：
\begin{enumerate}
    \item \textbf{时序分布失衡，短时动作难以召回}：
    手术视频具有显著的“长尾”数据特征，其流程通常由漫长的标准操作（如房角观察）和极其短暂的切换或精准交互操作组成。
    在这种分布极度不平衡的状态下，一些长阶段占据了数分钟时长，而决定阶段切换的关键动作（如注射粘弹剂）往往只有数秒。
    这导致模型极易产生预测偏向，倾向于维持高频的长阶段预测，
    而忽略低频但关键的短动作，导致关键切换节点的感知缺失，严重影响系统对手术状态判别的灵敏度。
    
    \item \textbf{外观高度相似，不同阶段难以区分}：
    在显微手术中，不同阶段的操作往往共用相同的手术器械，且背景环境几乎恒定（如始终聚焦于眼球且组织不会发生显著形变）。
    例如，“液体清洗”与“粘弹剂注入”在视觉特征上可能仅表现为细微的流体质感差异或者器械深入长短的差距。
    在在线识别场景下，模型无法像“事后复盘”那样查看后续画面来反向印证当前动作，只能依赖已有的画面信息。
    当不同阶段的视觉表现高度雷同，仅凭画面的时序特征很难实现精确区分，极易产生判断震荡。
\end{enumerate}


%% 现有工作不足：
% 现有的在线手术阶段识别研究主要侧重于构建复杂的时序模型（如 RNN、TCN 或 Transformer）以捕捉视频的长程依赖。
% 为满足实时性需求并解决数据分布问题，近期研究如 SKiT\cite{liu2023skit} 引入了高效的键聚类池化操作，
% 而 Meta-SurDiff\cite{li2025metasurdiff} 尝试利用分类扩散模型对预测分布进行细化。
% 尽管这些方法在处理通用手术流程方面取得了进展，但仍存在一定局限。

% 首先，针对“短时瞬态动作”的感知机制仍停留于静态语义层面。
% 现有的生成式或池化类方法本质上更倾向于对视频序列进行“全局分布拟合”，
% 即通过统计图像中器械与解剖结构的整体响应频率来定性，而忽略了帧间微弱的动态演化特征。
% 在青光眼手术中，诸如针尖刺入瞬间的特征突变往往被长时段的静态背景信息所淹没，
% 导致模型对关键动作切换点的感知缺乏物理意义上的“过程灵敏度”，难以在视觉干扰下实现鲁棒的特征增强。


%% 上下文切换管理：ProTAS（shen2024progress）（采用了Task Graph Learner），
%% SurgPlan++（chen2025surgplan++）（采用了FSM），ActionSwitch（注意针对overlap的设计）% 类似方案：ding2024neural（fsm）

% 其次，针对“视觉歧义”的手术规程约束缺乏环境适应性。
% 虽然 ProTAS\cite{shen2024protas} 与 NFSM\cite{ding2025nfsm} 分别尝试引入任务图与状态转移先验来强制维护时序连贯性，
% 但其机理主要建立在对训练集阶段跳转频率的“统计关联”建模上。
% 这种建模方式将规程逻辑视为一种硬性的后期修正，缺乏对实时场景中“环境置信度”与“内容逻辑”的深度解耦。
% 具体而言，模型无法区分当前的预测不确定性究竟是源于视频画面本身的质量降级（如术中出血或遮挡导致的感知失效），
% 还是来自动作序列与规程先验的逻辑冲突。
% 这种局限性导致模型在视觉清晰时可能受到逻辑先验的过度干扰，而在视觉模糊时又无法产生有效的“逻辑拉回”作用，
% 难以实现噪声外观特征与结构化规程知识之间的深度语义对齐。

% 综上所述，如何设计一种能感知瞬时动作动态差异的特征提取算子，
% 并构建一套能根据环境置信度自适应调整约束强度的逻辑校验引擎，是实现鲁棒在线识别的关键。


为解决上述挑战，本章在现有在线检测框架的基础上，提出了一种结合特征增强与逻辑修正的方法。
针对瞬时动作捕捉难题，我们设计了\textbf{自适应多尺度差分模块（A-MSDD）}。
该模块通过在线计算不同时距的特征差分，显式强化了瞬态变化信号，使模型能够精确锁定实时流中的阶段切换点。

针对语义混淆问题，本章进一步提出了\textbf{带门控的图逻辑插件（L-GBP）}。
该插件将手术逻辑转化为可进化的图表征，并配合双门控机制。
通过对当前环境置信度的实时评估，模型能自适应地利用逻辑先验（模拟医生经验）对视觉判别进行纠偏，显著增强了在线识别的一致性。

此外，本章提出了一种\textbf{自适应时序一致性损失函数（SA-Loss）}。
该损失函数能根据特征波动率动态调节监督力度，在保证在线识别低延迟的同时，有效平衡了切换边界的响应速度与阶段内部的平滑度。
通过上述模块的有机结合，本章构建了一个轻量、高效的在线手术阶段识别框架。


\section{问题定义}

在线动作检测（OAD）任务旨在从实时视频流中识别当前时刻正在发生的动作，并在每一帧即时预测其对应的类别标签。
给定一个持续输入的视频流 $V=\{f_t\}_{t=1}^T$，其中 $f_t$ 表示在时刻 $t$ 到达的视频帧，
任务的目标是根据历史观察及当前帧 $f_{1:t}$，在不依赖未来信息的前提下，
预测当前时刻 $t$ 的动作类别标签 $a_t \in \{0, 1, \dots, C\}$。在此定义中，
$C$ 表示预定义的动作类别总数，而标签 $0$ 通常被定义为“背景”类，
用于表示当前时刻没有任何目标动作发生。
不同于离线检测，在线动作检测要求模型在时间步 $t$ 实时输出预测结果 $\hat{a}_t$，即在每一个新的时间切片到达时，
建立从历史特征序列到当前动作状态的即时映射。

\section{本章方法}


\section{实验结果}

\subsection{数据集与评价指标}
\textbf{数据集构建}\quad

\textbf{评价指标}\quad

\textbf{实验细节}\quad

\subsection{性能比较}
% 与通用oad的性能比较（LSTR，TesTra, E2E-LOAD，CMerT，MerT）；
% 与在线手术识别模型的性能比较（SKViT，Meta-SurDiff，SurgPlan++）；


\subsection{消融分析}

\section{本章小结}
% \include{sec/5-plt}
\chapterx{结论}

% 在结论中应明确指出本研究内容的创造性成果或创新性理论（含新见解、新观点），
% 对其应用前景和社会、经济价值等加以预测和评价，并指出今后进一步在本研究方向进行研究工作的展望与设想。

\makeatletter
\defbibheading{mybib}[\bibname]{%
  % \cleardoublepage
  \phantomsection
  \bookmarksetup{startatroot}%
  \chapterxtrue
  \chapter*{#1}%
  \chapterxname{#1}%
  \addcontentsline{toc}{chapter}{#1}%
  \refbodyfont
  \let\ps@plain=\ps@revtitlestyle
  \let\ps@headings=\ps@revtitlestyle
}
\makeatother
\printbibliography[heading=mybib]
% % 不单独生成“附录”标题页，只在目录中登记“附录”
\clearpage
\phantomsection
\addcontentsline{toc}{chapter}{附~~录}
\chapterxtrue
\chapterxname{附~~录}

\renewcommand{\thetable}{附表\arabic{table}} 
\setcounter{table}{0} 
\renewcommand{\tablename}{} 
\captionsetup[table]{name=}

\includepdf[pages=1,pagecommand={\thispagestyle{revtitlestyle}\begin{center}{\appendixtitlefont 附~~录}\end{center}},scale=0.94,offset=0 -1.2cm]{\detokenize{assets/应用证明-中山三院.pdf}}
\includepdf[pages=-,pagecommand={\thispagestyle{revtitlestyle}},fitpaper=true]{\detokenize{assets/基于深度学习的微创青光眼手术实时导航系统应用证明.pdf}}

% \begin{table*}[h]
% \small
% \renewcommand{\arraystretch}{1.0}
% \renewcommand{\tabcolsep}{4pt}
% \begin{center}
% \caption{现有时序动作定位方法在各个微动作上的具体结果} \label{tab: detailed  tal}
% \begin{tabular}{llrrrrrrrr}
% \toprule
%  \multirow{2}{*}{方法}&\multirow{2}{*}{微动作}& \multicolumn{8}{c}{AP@tIoU } \\
%    && 0.1 & 0.2 & 0.3 & 0.4 & 0.5 & 0.6 & 0.7 & Avg.\\
% \midrule
%  \multirow{8}{*}{A2Net \cite{yang2020revisiting}}&口腔运送 & 42.3& 41.9& 34.1& 26.4& 13.1& 7.0& 2.3& 23.9\\ 
%   &软腭抬升  & 72.1& 66.8& 63.0& 48.8& 30.0& 9.6& 3.2& 41.9\\
%   &舌骨运动 & 60.0& 57.5& 53.3& 39.6& 25.6& 13.4& 4.5& 36.3\\ 
%   &UES开放 &  68.2& 38.8& 26.0& 13.7& 4.0& 0.9& 0.7& 17.9\\ 
%   &吞咽启动 & 12.0& 9.0& 3.9& 0.9& 0.3& 0.1& 0.1& 3.8\\ 
%   &咽期运送 & 76.8& 76.8& 75.4& 66.8& 48.4& 24.6& 7.8& 53.8\\ 
%   &喉前庭闭合  & 72.1& 69.8& 63.9& 55.8& 35.5& 18.6& 6.2& 46.0\\
% %  & \cell 平均& \cell53.7& \cell51.5& \cell45.7& \cell36.0& \cell22.4& \cell10.6& \cell3.5& \cell31.9\\ 
% \midrule
%  \multirow{8}{*}{ActionFormer \cite{zhang2022actionformer}}&口腔运送 & 48.2& 47.4& 40.9& 25.0& 14.5& 11.2& 7.6& 27.8\\ 
%   &软腭抬升  & 65.4& 64.6& 64.6& 57.5& 49.4& 36.9& 22.4& 51.6\\
%   &舌骨运动 & 73.9& 51.9& 35.2& 17.9& 9.4& 3.0& 1.3& 27.5\\ 
%   &UES开放 & 68.2& 68.1& 68.1& 67.0& 61.0& 52.6& 34.7& 60.0\\ 
%   &吞咽启动 & 35.7& 26.5& 21.1& 11.6& 6.7& 3.6& 2.2 & 15.4\\ 
%   &咽期运送 & 70.2& 70.2& 69.9& 54.7& 34.4& 23.5& 15.6& 48.4\\ 
%   &喉前庭闭合  & 83.4& 83.2& 82.9& 81.7& 71.8& 57.4& 39.1& 71.4\\
% %  & \cell 平均& \cell76.8& \cell74.4& \cell69.7& \cell59.3& \cell48.8& \cell38.5& \cell24.6& \cell56.0\\ 
% \midrule
%  \multirow{8}{*}{TriDet \cite{shi2023tridet}}& 口腔运送 & 45.8& 43.0& 38.5& 22.5& 15.8& 10.9& 6.1& 26.1\\ 
%   &软腭抬升  & 70.9& 70.9& 64.4& 60.3& 54.1& 41.7& 25.2& 55.4\\
%   &舌骨运动 & 73.8& 52.3& 41.3& 25.2& 17.1& 6.7& 2.7& 31.3\\ 
%   &UES开放 & 68.8& 68.7& 68.4& 66.1& 62.4& 51.3& 38.1& 60.5\\ 
%   &吞咽启动 & 39.0& 29.3& 20.9& 13.8& 8.2& 4.3& 2.1& 16.8\\ 
%   &咽期运送 & 71.6& 71.6& 67.9& 52.7& 37.9& 27.6& 16.9& 49.5\\ 
%   &喉前庭闭合  & 84.6& 84.4& 83.9& 83.4& 69.2& 58.0& 36.9& 71.5\\
% %  & \cell 平均& \cell79.6& \cell76.6& \cell72.4& \cell63.7& \cell53.1& \cell40.9& \cell26.2& \cell58.9\\ 
% \midrule
%  \multirow{8}{*}{AdaTAD \cite{liu2024end}}&口腔运送 & 73.4 & 69.9 & 46.3 & 36.1 & 24.9 & 11.2 & 4.6 & 38.0 \\ 
%   &软腭抬升  & 68.4 & 68.4 & 68.3 & 66.2 & 58.8 & 52.4 & 21.5 & 57.7 \\
%   &舌骨运动 & 81.4 & 62.8 & 35.8 & 22.2 & 17.0 & 9.5  & 5.7  & 33.5 \\ 
%   &UES开放 & 80.6 & 80.2 & 80.2 & 80.2 & 79.2 & 69.9 & 45.4 &  73.7 \\ 
%   &吞咽启动 & 29.8 & 20.7 & 12.3 & 9.3  & 6.6  & 3.8  & 1.6  & 12.0 \\ 
%   &咽期运送 & 65.0 & 65.0 & 64.8 & 57.3 & 48.7 & 29.9 & 16.4  & 49.6 \\ 
%   &喉前庭闭合  & 87.3 & 87.3 & 87.3 & 81.5 & 76.0 & 68.3 & 49.2  & 76.7 \\
% %  & \cell 平均& \cell80.9& \cell77.4& \cell70.0& \cell62.3& \cell54.4& \cell42.1& \cell24.8& \cell58.9\\ 
% \midrule
%  \multirow{8}{*}{ActionMamba \cite{chen2024video}}& 口腔运送 & 43.4 & 41.4 & 36.5 & 31.2 & 18.7 & 11.6 & 6.6 & 27.1\\ 
%   &软腭抬升 & 68.7 & 68.7 & 65.6 & 59.8 & 50.6 & 32.8 & 17.9 & 52.0 \\
%   &舌骨运动 & 73.1 & 50.2 & 37.1 & 26.4 & 15.5 & 8.2  & 3.3 & 30.5\\ 
%   &UES开放  & 69.0 & 68.8 & 68.8 & 68.3 & 63.8 & 45.1 & 27.4 & 58.8\\ 
%   &吞咽启动 & 46.5 & 36.8 & 26.4 & 17.6 & 10.1 & 6.2 & 2.2 &  20.8 \\ 
%   &咽期运送 & 77.6 & 77.6 & 69.8 & 48.3 & 29.0 & 20.6 & 10.6 & 47.6 \\ 
%   &喉前庭闭合  & 82.9 & 82.0 & 82.0 & 82.0 & 72.6 & 57.1 & 33.6 &  70.3 \\ 
% %  & \cell 平均&  \cell82.0	& \cell79.3& \cell74.3 &\cell67.4 & \cell53.1	&\cell37.1&	\cell19.9&	\cell59.0 \\
%  \midrule
%  \multirow{8}{*}{\sexyname \\（本文方法）}  &口腔运送  & 73.6 & 69.7 & 62.7 & 51.1 & 39.5 & 21.5 & 12.2 & 47.2 \\ 
%   &软腭抬升  & 86.4 & 84.7 & 84.6 & 82.9 & 74.6 & 67.8 & 46.2 & 75.3 \\
%   &舌骨运动 & 88.0 & 87.5 & 82.6 & 67.1 & 37.0 & 23.4 & 6.6 & 56.0 \\ 
%   &UES开放 & 85.9 & 85.9 & 85.5 & 85.2 & 85.0 & 83.6 & 77.8 & 84.1 \\ 
%   &吞咽启动 & 73.0 & 47.9 & 34.0 & 18.0 & 10.9 & 5.1 & 2.2 & 27.3 \\ 
%   &咽期运送 & 90.7 & 89.8 & 89.3 & 89.3 & 86.5 & 75.2 & 54.6 & 82.2 \\ 
%   &喉前庭闭合  & 84.2 & 84.0 & 83.4 & 82.1 & 79.3 & 73.5 & 60.9 & 78.2 \\
% %   & \cell 平均& \cell83.1& \cell78.5& \cell74.6& \cell67.9& \cell59.0& \cell50.0& \cell37.2& \cell64.3\\ 
%  \bottomrule
% \end{tabular}

% \label{tab: actions}
% \end{center}
% \end{table*}

% \begin{table*}[h]
% \small
% \renewcommand{\arraystretch}{1.0}
% \renewcommand{\tabcolsep}{3pt}
% \begin{center}
% \caption{更多方法在各个微动作上的详细结果比较} \label{tab: detailed  main}
% \begin{tabular}{lcclrrrrrrrr}
% \toprule
%  \multirow{2}{*}{方法}&\multirow{2}{*}{
% 检测器} &\multirow{2}{*}{骨架}&\multirow{2}{*}{微动作}& \multicolumn{8}{c}{AP@tIoU } \\
%    & &&& 0.1 & 0.2 & 0.3 & 0.4 & 0.5 & 0.6 & 0.7 & Avg.\\
% \midrule
%  \multirow{8}{*}{Ruan \textit{et al.} \cite{ruan2023temporal}}&\multirow{8}{*}{A2Net} &&口腔运送 & 75.1& 72.2& 59.8& 38.7& 23.0& 11.4& 7.1& 41.0\\ 
%   & &&软腭抬升  & 71.7& 71.7& 71.7& 71.2& 66.3& 42.9& 10.1& 57.9\\
%   & &&舌骨运动 & 75.4& 75.4& 71.8& 50.9& 27.9& 7.8& 1.1& 44.3\\ 
%   & &&UES开放 & 71.7& 71.7& 71.7& 71.7& 71.2& 57.5& 42.1& 65.4\\ 
%   & &&吞咽启动 & 58.7& 37.7& 19.3& 8.8& 3.8& 1.5& 8.0& 19.7\\ 
%   & &&咽期运送 & 71.9& 71.9& 71.9& 71.9& 65.0& 43.7& 19.8& 59.4\\ 
%   & &&喉前庭闭合  & 71.7& 71.7& 71.7& 71.7& 65.6& 56.6& 29.4& 62.6\\
% %  &  && \cell 平均& \cell70.9& \cell67.5& \cell62.5& \cell55.0& \cell46.1& \cell31.6& \cell15.8& \cell49.9\\ 
% \midrule
%  \multirow{8}{*}{Ruan \textit{et al.} \cite{ruan2023temporal}}&\multirow{8}{*}{ActionMamba} &&口腔运送 & 70.9& 68.8& 57.4& 42.4& 37.7& 27.6& 14.2& 45.6\\ 
%   & &&软腭抬升  & 80.8& 80.7& 78.4& 78.1& 64.9& 42.7& 20.4& 63.7\\
%   & &&舌骨运动 & 74.4& 73.8& 69.6& 56.6& 34.7& 21.3& 7.4& 48.2\\ 
%   & &&UES开放 & 80.5& 80.5& 79.7& 79.7& 79.1& 76.1& 63.8& 77.1\\ 
%   & &&吞咽启动 & 71.4& 55.4& 37.3& 17.6& 12.9& 6.7& 2.2 & 29.1\\ 
%   & &&咽期运送 & 84.2& 83.7& 82.2& 80.4& 76.8& 71.6& 41.7& 74.4\\ 
%   & &&喉前庭闭合  & 83.2& 82.1& 82.0& 80.7& 77.9& 70.5& 52.6& 75.6\\
% %  &  && \cell 平均& \cell77.9& \cell75.0& \cell69.5& \cell62.2& \cell54.8& \cell45.2& \cell28.9& \cell59.1\\ 
% \midrule
%  \multirow{8}{*}{Hyder \textit{et al.} \cite{hyder2024action}}&\multirow{8}{*}{ActionMamba}  &\multirow{8}{*}{\checkmark}& 口腔运送 & 71.7& 69.6& 60.1& 46.2& 40.7& 29.1& 13.7& 47.3\\ 
%   & &&软腭抬升  & 83.6& 83.6& 80.0& 78.0& 72.7& 64.8& 46.6& 72.3\\
%   & &&舌骨运动 & 81.4& 78.9& 76.9& 54.6& 31.4& 13.9& 1.6& 48.4\\ 
%   & &&UES开放 & 80.6& 80.4& 80.0& 79.8& 79.2& 71.0& 44.2& 73.6\\ 
%   & &&吞咽启动 & 53.6& 43.1& 24.4& 12.8& 7.7& 3.7& 2.0& 21.0\\ 
%   & &&咽期运送 & 86.0& 86.0& 85.5& 84.3& 83.4& 68.8& 52.4& 78.0\\ 
%   & &&喉前庭闭合  & 80.6& 80.6& 79.2& 78.6& 78.3& 68.4& 54.9& 74.4\\
% %  &  && \cell 平均& \cell76.8& \cell74.6& \cell69.4& \cell62.0& \cell56.2& \cell45.2& \cell30.8& \cell59.3\\ 
%  \midrule
%  \multirow{8}{*}{\sexyname \\（本文方法）}&\multirow{8}{*}{ActionFormer} &\multirow{8}{*}{\checkmark}&
%       口腔运送 & 75.4 & 71.2 & 66.7 & 47.9 & 34.3 & 19.4 & 11.0 & 46.6 \\
% & &&软腭抬升 & 86.6 & 85.1 & 83.8 & 82.5 & 75.4 & 66.4 & 42.2 & 74.6 \\
% & &&舌骨运动 & 88.2 & 87.3 & 82.2 & 65.6 & 42.2 & 26.8 & 3.6 & 56.5 \\
% & &&UES开放 & 85.8 & 85.8 & 85.3 & 84.8 & 84.1 & 81.9 & 74.0 & 83.1 \\
% & &&吞咽启动  & 66.6 & 56.0 & 40.0 & 16.7 & 10.5 & 5.8 & 2.6 & 28.3 \\
% & &&咽期运送 & 92.4 & 91.6 & 89.8 & 88.8 & 84.1 & 70.8 & 48.2 & 80.8 \\
% & &&喉前庭闭合 & 83.8 & 83.8 & 83.2 & 83.2 & 78.3 & 73.5 & 53.1 & 77.0 \\
% % & &&\cell 平均 & \cell 82.7 & \cell 80.1 & \cell 75.9 & \cell 67.1 & \cell 58.4 & \cell 49.2 & \cell 33.5 & \cell 63.8 \\ 
% \midrule
%  \multirow{8}{*}{\sexyname \\（本文方法）} &\multirow{8}{*}{ActionMamba} &\multirow{8}{*}{\checkmark}&口腔运送 & 73.6 & 69.7 & 62.7 & 51.1 & 39.5 & 21.5 & 12.2 & 47.2 \\ 
%   & &&软腭抬升  & 86.4 & 84.7 & 84.6 & 82.9 & 74.6 & 67.8 & 46.2 & 75.3 \\
%   & &&舌骨运动 & 88.0 & 87.5 & 82.6 & 67.1 & 37.0 & 23.4 & 6.6 & 56.0 \\ 
%   & &&UES开放 & 85.9 & 85.9 & 85.5 & 85.2 & 85.0 & 83.6 & 77.8 & 84.1 \\ 
%   & &&吞咽启动 & 73.0 & 47.9 & 34.0 & 18.0 & 10.9 & 5.1 & 2.2 & 27.3 \\ 
%   & &&咽期运送 & 90.7 & 89.8 & 89.3 & 89.3 & 86.5 & 75.2 & 54.6 & 82.2 \\ 
%   & &&喉前庭闭合  & 84.2 & 84.0 & 83.4 & 82.1 & 79.3 & 73.5 & 60.9 & 78.2 \\
% %  &  && \cell 平均& \cell83.1& \cell78.5& \cell74.6& \cell67.9& \cell59.0& \cell50.0& \cell37.2& \cell64.3\\ 
%  % \midrule
%  % \multirow{8}{*}{\sexyname (Oracle)}&\multirow{8}{*}{ActionMamba} &\multirow{8}{*}{\checkmark}&口腔运送 & 91.3& 89.4& 87.2& 61.4& 41.1& 22.5& 12.0& 57.8\\ 
%  %  & &&软腭抬升  & 92.2& 91.2& 91.1& 91.1& 85.1& 76.1& 49.6& 82.3\\
%  %  & &&舌骨运动 & 94.5& 93.8& 91.7& 78.4& 58.0& 33.4& 10.2& 65.7\\ 
%  %  & &&UES开放 & 93.2& 93.2& 93.2& 93.2& 92.1& 91.4& 88.8& 92.1\\ 
%  %  & &&吞咽启动 & 80.2& 49.7& 35.4& 19.6& 11.0& 5.5& 2.6& 29.2\\ 
%  %  & &&咽期运送 & 93.6& 93.6& 93.6& 93.6& 92.0& 85.6& 58.6& 87.2\\ 
%  %  & &&喉前庭闭合  & 93.6& 93.6& 93.6& 92.6& 88.8& 82.5& 73.9& 88.4\\
%  % &  && \cell 平均& \cell{91.2} &\cell{86.3} &	\cell{83.7}& 	\cell{75.7} &	\cell{66.9}& 	\cell{56.7}& 	\cell{42.2} 	&\cell{71.8} \\ 
%  \bottomrule
% \end{tabular}
% \label{tab:detailed actions}
% \end{center}
% \end{table*}


% \begin{algorithm*}[h]
% \caption{\crossblockfull(\crossblockshort)模块的PyTorch风格的伪代码实现}
% \label{alg: pytorch ccm}
% \begin{minted}
% [fontsize=\footnotesize, baselinestretch=1]
% {python}
% def forward(self, feat, com_feat):
%     """
%     输入:
%         feat:      (B, L, D)     - 输入特征
%         com_feat:  (B, L, D')    - 互补特征
%     输出:
%         enhanced_feat: (B, L, D) - 增强特征
%     """
%     # 1. 特征投影
%     x = self.in_proj(feat)          # (B, L, D) -> (B, 2D, L)
%     z = self.v_in_proj(com_feat)    # (B, L, D') -> (B, 2D, L)

%     # 2. 双向分块
%     x_f, x_b = torch.chunk(x, 2, dim=1)  # (B, 2D, L) -> [(B, D, L), (B, D, L)]
%     z_f, z_b = torch.chunk(z, 2, dim=1)  # (B, 2D, L) -> [(B, D, L), (B, D, L)]
%     x_b = x_b.flip([-1])  # Reverse the backward chunk
%     z_b = z_b.flip([-1])  # Reverse the backward chunk

%     # 3. 拼接特征
%     x = torch.cat([x_f, x_b], dim=0)  # (2B, D, L)
%     z = torch.cat([z_f, z_b], dim=0)  # (2B, D, L)
%     xz = torch.cat([x, z], dim=1)     # (2B, 2D, L)

%     # 4. SSM交互
%     out = self.conv_and_ssm(xz)  # (2B, 2D, L) -> (2B, 2D, L)

%     # 5. 恢复顺序
%     out = out.chunk(2)  # Split into forward and backward outputs
%     out = torch.cat([out[0], out[1].flip([-1])], dim=1)  # (B, 2D, L)

%     # 6. 通道增强
%     comm = F.silu(self.conv(out))   # (B, 2D, L) -> (B, 1, L)
%     out = out + self.gamma * (out - comm)

%     # 7. 特征融合
%     enhanced_feat = self.out_proj(out) + feat  # (B, 2D, L) -> (B, L, D)

%     return enhanced_feat
% \end{minted}
% \end{algorithm*}


\chapterx{攻读硕士学位期间取得的研究成果}

一、已发表（包括已接受待发表）的论文，以及已投稿、或已成文打算投稿、或拟成文投稿的论文情况（只填写与学位论文内容相关的部分）：

\begin{longtable}{|>{\centering}m{0.5cm}|>{\centering}m{3.5cm}|>{\centering}m{2.3cm}|>{\centering}m{2.6cm}|>{\centering}m{2cm}|>{\centering}m{1.3cm}|>{\centering}m{0.9cm}|}
\hline 
序号 & 作者（全体作者，按顺序排列） & 题 目 & 发表或投稿刊物名称、级别 & 发表的卷期、年月、页码 & 相当于学位论文的哪一部分（章、节） & 被索引收录情况\tabularnewline
\hline 
1 & \textbf{李奕锐}，周凯，戴萌，周海宇，胡晋武，杨逸凡，陈健，刘飞，蔡宏民，谭明奎 & Temporal Micro-Action Localization with Skeleton-Guided Mamba for Videofluoroscopic Swallowing Study & IEEE Transactions on Neural Systems and Rehabilitation Engineering & 已投稿 & 第三章 & - \tabularnewline
\hline 
2 & 张寅航，\textbf{李奕锐}（共同一作），周凯，
， 汪德明， 杨泽锋， 游增， 陈瑜， 黎晓燕， 闫晓伟， 汪之璇， 蒋嘉烜， 
金玲， 王贞玉， 唐广贤， 范肃洁， 李飞， 谭明奎， 张秀兰
& Multicenter fine annotated surgical video dataset for minimally invasive glaucoma surgery & Scientific Data & 已投稿 & 第四章 & - \tabularnewline
\hline
\end{longtable}

注：在“发表的卷期、年月、页码”栏：
1如果论文已发表，请填写发表的卷期、年月、页码；
2如果论文已被接受，填写将要发表的卷期、年月；
3以上都不是，请据实填写“已投稿”，“拟投稿”。

二、与学位内容相关的其它成果（包括专利、著作、获奖项目等）

专利：已授权一项发明专利《一种吞咽造影视频微动作识别与定位方法及电子设备》，第四发明人，2025

\quad\quad 已受理一项发明专利《一种基于深度学习的青光眼手术实时导航方法》，第五发明人，2025

% \chapterx{攻读硕士学位期间取得的研究成果}

一、已发表（包括已接受待发表）的论文，以及已投稿、或已成文打算投稿、或拟成文投稿的论文情况\textbf{\underline{（只填写与学位论文内容相关的部分）：}}

\begin{longtable}{|>{\centering}m{0.5cm}|>{\centering}m{3.5cm}|>{\centering}m{2.3cm}|>{\centering}m{2cm}|>{\centering}m{2.6cm}|>{\centering}m{2cm}|}
\hline 
\textbf{序号} & \textbf{发表或投稿刊物/会议级别} & \textbf{作者（仅注明第几作者）}  & \textbf{发表年份} & \textbf{与学位论文的哪一部分（章、节）相关} & \textbf{被索引收录情况}\tabularnewline
\hline 
1 & IEEE Journal of Biomedical and Health Informatics & 第一作者 & 拟投稿  & 第三章 & - \tabularnewline
\hline 
2 & Scientific Data & 共同一作 &  已投稿 & 第四章 & - \tabularnewline
\hline
\end{longtable}

\noindent 注：1.请在“作者”一栏填写本人是第几作者，例：“第一作者”或“导师第一，本人第二”等；

\noindent\quad\quad2.若文章未发表或未被接受，请在“发表年份”一栏据实填写“已投稿”，“拟投稿”。

\noindent\quad\quad 不够请另加页。

二、与学位内容相关的其它成果（包括专利、著作、获奖项目等）

专利：已授权一项发明专利，第四发明人，2025

\quad\quad 已受理一项发明专利，第五发明人，2025

\chapterx{致~~谢}

时光荏苒，硕士阶段的学习与研究即将落幕。
回首这段充实而珍贵的求学之路，心中满是感恩与不舍，
谨在此向所有给予我帮助与支持的师长、同窗与家人挚友，致以最诚挚的谢意。

首先，我由衷感谢谭明奎教授。
谭老师以严谨的治学态度、深厚的学术素养和开阔的科研视野，
为我指明方向、解惑答疑，教会我踏实做事、认真治学。
谭老师的悉心教诲与耐心指引，是我顺利完成学业的坚实保障，让我受益终身。

感谢杜卿副教授在科研实践与论文完善过程中的悉心指导与无私帮助。
杜老师严谨细致、求真务实，让我的研究工作不断完善与精进。

感谢陈健教授在学术交流与课题研讨中的宝贵指导。
陈教授深入浅出的讲解与中肯的建议，拓宽了我的研究思路，使我对课题的理解与思考更为深入全面。

感谢实验室周凯、游增、胡晋武、杨一帆、张书海、
阮湘辉等诸位师兄的倾囊相授、耐心疑惑，
在实验、数据处理与论文写作上给予我大量帮助，
让我快速融入科研环境、稳步前行。
感谢向添航、罗海林、杨伟好、陈传深、康海龙、樊琚岗等同门的帮助和鼓励。

衷心感谢中山大学附属第三医院与中山大学中山眼科中心团队。
本研究的顺利完成，离不开医院团队提供的宝贵临床数据支撑与专业指导，
感谢各位老师与师兄师姐等在数据采集与合作研究中的辛勤付出与大力支持。

最后，感谢我的家人与亲密伙伴。你们无条件的理解、包容与默默支持，是我克服困难、坚持前行的最大底气与动力。
道阻且长，行则将至。未来我将带着这份感恩与期许，不忘初心、脚踏实地，在科研与人生道路上继续砥砺前行，不负所有师长、朋友与亲人的期望。

% \begin{minipage}[t]{0.8\columnwidth}%
% \begin{flushright}
% 李奕锐 
% \par\end{flushright}

% \begin{flushright}
% 2026年3月28日
% \par\end{flushright}%
% \end{minipage}

\end{document}
